\begin{multicols}{2}[\section*{Objetivos}]
\subsubsection*{Objetivo general}
\begin{itemize}[leftmargin=*, itemsep=2pt]
  \item Estimar la densidad lineal $\mu$ de cuerdas de algodón, goma delgada y goma gruesa usando ondas estacionarias.
\end{itemize}

\subsubsection*{Objetivos específicos}
\begin{enumerate}[leftmargin=*, itemsep=2pt, label=\textbf{\arabic*.}]
  \item Aplicar dos protocolos: \textbf{tensión variable} (con $\nu$ fija) y \textbf{frecuencia variable} (con $T$ fija), asignando \emph{una muestra distinta por protocolo} para cada material (total: seis cuerdas).
  \item Para cada muestra, extraer $\mu$ a partir de la pendiente del ajuste lineal correspondiente y contrastarla con la estimación geométrica $\mu_{0}=m/L$ de esa misma muestra.
  \item Discutir limitaciones del modelo ideal (fricción en polea, extensibilidad, definición de $L$ efectivo) y su impacto en las estimaciones.
\end{enumerate}
\end{multicols}


